\sectionguard
\section*{The Propers: Interpretive Possibilities}

\noindent\emph{These four connections are exploratory editorial and
AI-assisted proposals, not sourced historical intent or teaching attributed
to the witnesses cited elsewhere.}

\begin{proposal}{Burden-bearing and the halted bier\enspace\properrefs{Ep., Gosp.}}

\pfield{Anchors}{Galatians~6:2,5 pairs \latin{ónera portáte} and
\latin{onus suum portábit}; Luke~7:14 names the bearers,
\latin{qui portábant}, who halt when Christ touches the bier.}

\pfield{Mechanism}{The carrying verbs invite comparison of what help can
accomplish. Companions bear another's weight and accompany need; Christ
gives the life they cannot produce. Their pause before his command can
illuminate charity that remains active while acknowledging the divine gift
upon which restoration depends.}

\pfield{What separate treatment misses}{An account of moral burdens and
an account of a funeral can each be accurate without asking how helpful
action meets a need beyond its power. Their conjunction sets the capacity
for assistance beside its limit.}

\pfield{Fruit}{Responsibility can be energetic without becoming
self-sufficient. The helper may remain present, bear a real weight and
intercede while entrusting the outcome to Christ. Aquinas's connection
between burden-bearing and Christ's bearing of sins supplies a near
analogue (\emph{Super Galatas} 6.1), without this specific bier comparison.}

\pfield{Strongest limit}{Burial is a proper duty, and Luke does not
condemn the bearers. Bodily death is not identified with personal guilt;
accompaniment is not useless, and mutual aid does not dissolve the
personal account that Paul retains.}

\end{proposal}
\begin{proposal}{Unreported speech and a received song\enspace\properrefs{Gosp., Off.}}

\pfield{Anchors}{Luke~7:15 reports \latin{cœpit loqui} without quoting
the youth; the Offertory gives \latin{immísit in os meum cánticum novum}.}

\pfield{Mechanism}{Restored capacity for speech may be set beside the
assembly's later first-person thanksgiving. Worshippers can exercise a
voice received as gift without filling the narrator's silence with an
invented sentence. The relation concerns their participation in praise,
not recovery of the boy's unreported words.}

\pfield{What separate treatment misses}{The miracle makes speech evidence
of life; the psalm makes thanksgiving a received gift. Their conjunction
asks how received life becomes an act the worshipping community performs.}

\pfield{Fruit}{Thanksgiving can answer restoration even when another
person's experience remains beyond one's knowledge. \emph{The Liturgical
Year}, p.~354, already joins Mother Church's restored children to
Offertory praise; the further proposal considers the worshipper's voice
beside the narrative's unspecified speech.}

\pfield{Strongest limit}{Luke distinguishes the boy's speech from the
crowd's glorification. He narrates neither the boy singing nor his use of
Psalm~39. The psalm's speaker has an independent canonical and received
identity and cannot be replaced by the youth.}

\end{proposal}

\begin{proposal}{Opportunity within endurance\enspace\properrefs{Int., Grad., Ep., Off.}}

\pfield{Anchors}{The Introit's \latin{tota die}, the Gradual's
\latin{mane} and \latin{per noctem}, the Epistle's \latin{témpore enim suo}
and \latin{dum tempus habémus}, and the Offertory's
\latin{Exspéctans exspectávi} name recurrence, opportunity and expectation.}

\pfield{Mechanism}{Their several scales of time may form one practical
discipline. All-day petition and morning/night praise recur; an
opportunity for good presses now; the harvest awaits its fitting time.
Patient thanksgiving places both present action and delayed fruit within
dependence upon God.}

\pfield{What separate treatment misses}{Each expression can be explained
alone while the tension between urgent service and patient expectation
remains unexamined. Their conjunction asks how a person can hurry to help
without demanding an immediate result.}

\pfield{Fruit}{Perseverance may sustain prompt service through an interval
whose length the helper does not control. Aquinas's patient harvest
(\emph{Super Galatas} 6.2) and Schuster's target treatment
(III, pp.~139--141) offer temporal analogues to this fuller coordination.}

\pfield{Strongest limit}{These expressions form no hidden calendar or
successive series of liturgical phases. They promise no desired result on
schedule. Morning/night has different received interpretations, and
Psalm~39 returns from thanksgiving to need.}

\end{proposal}

\begin{proposal}{Heard prayer, seen grief, touch and public voice\enspace\properrefs{Int., Gosp., Off.}}

\pfield{Anchors}{The Introit names \latin{aurem} and \latin{exáudi};
Luke~7:13--16 joins seeing, touching, speaking and glorifying; the
Offertory returns to hearing and the mouth's canticle.}

\pfield{Mechanism}{Need and response become perceptible in different
forms. The servant requests a hearing, while the widow's grief is seen
before a request is narrated. Touch interrupts the procession; restored
speech and shared praise become audible. These acts suggest attention
capable of receiving more than a verbally stated request.}

\pfield{What separate treatment misses}{The earlier speech proposal asks
how worshippers take up thanksgiving beside unreported words. This one
asks how mercy attends to need that is not only spoken: prayer, perception
and concrete assistance enter the same encounter.}

\pfield{Fruit}{Compassion may begin by noticing a person's situation,
continue through appropriate action, and leave room for that person's
own voice and the community's thanksgiving. The relation joins embodied
attention to prayer without reducing either to the other.}

\pfield{Strongest limit}{The divine ear is not human anatomy. Jesus
touches the bier, not the corpse. Hearing, seeing, touch and speech form
no established developmental ladder, single speaker, or demonstrated
program of the formulary's compiler.}

\end{proposal}
